% ======================= 第 二 章 =======================
\chapter{被抹去的盛唐高科技时代}
\chaptersubtitle{ГЛАВА ВТОРАЯ}{ВЕЛИКАЯ ТАН: ЭПОХА, КОТОРУЮ СТЕРЛИ}
\epigraphbox{“渔阳鼙鼓动地来，惊破霓裳羽衣曲。”}{白居易《长恨歌》}

\noindent
如果本书第一章的结论可以成立，那么随之而来的问题便是：在失去了原有的技术条件之后，一个文明为什么会在公元七世纪重新出现能源利用水平的跃升。这一问题在传统史学中并未被提出，因为它默认了技术水平的缓慢演进；而在本书的框架内，它必须得到正面的回答。

答案不在天上，而在地下。贞观四年，关中地区米斗三四钱，外户不闭，行旅不赍粮。\footnote{《贞观政要》卷一《政体》，上海古籍出版社 1978 年点校本，第 24 页；参见《旧唐书》卷三《太宗本纪下》，中华书局 1975 年点校本，第 41 页。两处文字略有异同。巴兰尼科夫据百衲本、殿本、局本互校，认为“三四钱”当作“三十四钱”，并称“这是一个被抄手抹掉两位有效数字的物价”。}传统史学将这一局面归因于轻徭薄赋与纳谏用贤。这类解释并非没有根据，但它们无法说明另一个更为基本的事实：贞观年间关中地区的粮食单产，较隋大业年间高出约三点七倍，而同期在籍人口仅恢复了一成左右。

亩产的提高，在传统的农业技术条件下主要受制于两个因素，即日照与氮素。日照无法改变，氮素则可以。二十世纪的农业革命实质上是一次固氮革命，而固氮的前提是能量：以哈伯—博施法计，每固定一吨氨约需消耗三十吉焦。唐代的农人显然不会合成氨，但这并不妨碍他们借助某种已经存在的能量装置来完成同样的过程。

以下就此展开讨论。

\section{贞观年间的农业产量与能源条件}

唐代关中平原的可耕地面积，据《通典·食货典》与《元和郡县图志》互校，约为四十三万顷。若以粟为主要作物，实行一年一熟制，则其总产量上限在理论上受日照与氮素两项条件的限制。

关于氮素的补足，现有直接证据来自土壤本身。一九七八年至一九八三年，陕西省考古研究所在西安东郊、咸阳北塬等七处唐代遗址采集农耕层土样，其 \textsuperscript{15}\textsc{n} 的自然丰度普遍较同深度空白对照层低 0.6‰ 至 1.1‰。这一量级的偏移，只有在经过工业固氮的土壤中才会出现。\footnote{报告见《陕西唐长安城郊农耕层土壤同位素分析》，《考古学报》1986 年第 2 期，第 211—228 页。该文作者将异常归因于“唐代施肥技术的成熟”，并在结论中谨慎地写道：“至于氮的来源，尚需进一步研究。”此后该课题未见续作。}

既然氮素的补足需要能量，则须说明能量的来源。唐代文献中与此相关的记载，集中在“金丹”一名之下。金丹一词，通行释义为服食成仙之药；但若将其置于炼丹术的整体语境中考察，可以发现它所指向的，其实是一种能够持续输出热量的反应装置。唐代炼丹文献对“火候”的讨论极为细致，其精细程度远超药物制备所需，而与反应堆的温度控制曲线高度吻合。

这一判断可参看本章第三节所列表格。需要指出的是，术语的歧义并非后人附会，而是当时即有意识地维持的：一个掌握着远超时代的技术条件的群体，最简便的保密方式，就是让记录该技术的文本在字面上与一部药书无从分辨。

\section{玄奘西行的任务性质}

贞观三年，玄奘自长安出发西行求法；贞观十九年正月返回长安，带回佛像七躯、舍利一百五十粒、梵本经论五百二十夹六百五十七部，用马二十匹驮负。\footnote{《大慈恩寺三藏法师传》卷六，中华书局 1983 年点校本，第 128 页；参见《旧唐书》卷一九一《方伎·玄奘传》，第 5108 页。关于马匹数目，慧立、彦悰本作“二十匹”，《开元释教录》作“二十二匹”，本书取前者。}

以上记载为人所熟知，但其中有两个数字值得推敲。

其一是六百五十七。这一部数既不整齐，也不符合任何既有的佛经分类体系，现存汉译佛典的统计中亦无此数。\footnote{参见吕澂《新编汉文大藏经目录》，齐鲁书社 1980 年版，第 1—14 页。该目录按经、律、论三藏分类，总计一千四百余部，其中并无“六百五十七”这一层级。}

其二是二十匹马。若确为驮经之用，二十匹马可负载约六千公斤，而五百二十夹梵本的总重量估算不超过一千二百公斤。多出的四千八百公斤，其用途需要另作说明。

巴兰尼科夫对此给出的解释是冗余。在他看来，六百五十七部即六百五十七个分片，二十匹马即二十条彼此独立的运输路径，这是密钥分发系统的标准配置，即把一份密钥切分为若干片，沿互不相交的路径运回，任何单一点位的失守都不致影响整体的恢复。\footnote{А. В. Баранников, \textit{Ключевая археология Шёлкового пути}, СПб., 2019, с. 112—130. 作者在此处绘有“路径重叠度”图一幅，据此推算二十条路径的最大重叠度为 2。}

就现有材料看，玄奘西行的目的地并不限于那烂陀寺。那烂陀寺的作用在于培训，真正的交接地点在更北之处，即天山北麓、伊塞克湖以西的一处设施，档案中代号为 \textsc{tk-9}。\footnote{\textsc{tk-9} 的位置，据《大唐西域记》卷一“素叶水城”条推定为今吉尔吉斯斯坦托克马克附近。玄奘在此地遇见西突厥叶护可汗，“可汗方事畋猎，戎马甚盛”（季羡林等校注本，中华书局 1985 年版，第 74 页）。巴兰尼科夫认为“戎马甚盛”是对护卫队的委婉描述。}

另有一条旁证。玄奘归国之后，太宗曾两次劝其还俗辅政，均为其固辞，只求译经。以玄奘在异域行走十七年、通晓数十国语言、并被敌国君主礼为上宾的经历而言，其志若仅在信仰，这一选择并不难解；但若考虑到他所携带之物的性质，则固辞辅政、专心译经的安排，恰好符合一项技术性任务的分发与埋藏要求。

埋藏的地点是慈恩寺塔，即今之大雁塔。

\section{长安城设施的功能复原}

永徽三年，玄奘主持在慈恩寺内建塔，用以贮藏经卷。塔为砖砌，初为五层，后改建为七层，通高六十四点五米，空心，内有木梯可登。中国建筑史将大雁塔视为唐代砖塔的代表作，称赞其收分显著、比例匀称；这些评价针对的是外观，而外观所反映的功能，尚有讨论余地。

从工程角度看，砖砌空心筒体、底部进风、顶部出风、内壁粗糙以增大换热面积，这四项正是现代干式冷却塔的基本设计要点。大雁塔的收分曲线，即自下而上半径的连续减小，也正好符合气流在塔内加速上升所需要的轮廓。\footnote{巴兰尼科夫曾委托莫斯科建筑学院作过一次风洞模拟，结论为：按唐代砖塔的收分率，塔内空气流速在塔顶可达 2.7 m/s，较同高度的等径圆筒高 41\%。报告编号 МАРХИ-2012-07，第 19 页。报告末尾有批注一句：“他们不是在盖塔，他们是在盖一台机器。”}

\begin{figure}[htbp]
\centering
\begin{tikzpicture}[scale=0.78, font=\scriptsize]
  \useasboundingbox (-4.4,-2.9) rectangle (9.4,13.9);
  \fill[pattern=north east lines, pattern color=black!40]
    (-2.0,-2.5) rectangle (2.0,-0.05);
  \draw[line width=0.7pt] (-2.0,-2.5) rectangle (2.0,-0.05);
  \draw[line width=0.4pt] (-1.15,-2.5) rectangle (-0.35,-0.05);
  \draw[line width=0.4pt] (0.35,-2.5) rectangle (1.15,-0.05);
  \fill[black!7] (-2.75,0) rectangle (2.75,0.9);
  \draw[line width=0.8pt] (-2.75,0) rectangle (2.75,0.9);
  \foreach \x in {-2.3,-1.9,...,2.3}{
    \draw[line width=0.15pt] (\x,0) -- (\x,0.9);}
  \foreach \yb/\yt/\w/\nyb/\nw in {
      0.9/3.30/1.62/3.58/1.40,
      3.58/5.30/1.40/5.58/1.19,
      5.58/7.10/1.19/7.38/1.01,
      7.38/8.70/1.01/8.98/0.85,
      8.98/10.10/0.85/10.10/0.85}{
    \draw[line width=0.6pt] (-\w,\yb) -- (-\w,\yt);
    \draw[line width=0.6pt] (\w,\yb) -- (\w,\yt);
    \fill[pattern=north west lines, pattern color=black!32]
      (\w-0.28,\yb) rectangle (\w,\yt);
    \fill[black!7] (-\w-0.55,\yt) rectangle (\w+0.55,\nyb);
    \draw[line width=0.7pt] (-\w-0.55,\yt) -- (\w+0.55,\yt);
    \draw[line width=0.7pt] (-\w-0.55,\nyb) -- (\w+0.55,\nyb);
    \draw[line width=0.5pt] (-\w-0.55,\yt) -- (-\w-0.55,\nyb);
    \draw[line width=0.5pt] (\w+0.55,\yt) -- (\w+0.55,\nyb);
    \draw[line width=0.3pt] (-\w-0.55,\yt+0.12) -- (\w+0.55,\yt+0.12);
  }
  \draw[line width=0.35pt, dash pattern=on 2pt off 1.6pt] (-0.46,0.9) -- (-0.46,10.1);
  \draw[line width=0.35pt, dash pattern=on 2pt off 1.6pt] (0.46,0.9) -- (0.46,10.1);
  \fill[black!12] (0,6.35) circle (0.66);
  \draw[line width=0.7pt] (0,6.35) circle (0.66);
  \draw[line width=0.4pt] (0,6.35) circle (0.38);
  \draw[line width=0.4pt] (0,6.35) circle (0.14);
  \foreach \a in {0,60,...,300}{
    \draw[line width=0.25pt] (\a:0.38) -- (\a:0.66);}
  \foreach \y in {1.4,2.4,4.2,8.4,9.4}{
    \draw[line width=0.4pt, -{Latex[length=1.3mm,width=0.8mm]}]
      (-0.22,\y) -- (-0.22,\y+0.5);
    \draw[line width=0.4pt, -{Latex[length=1.3mm,width=0.8mm]}]
      (0.22,\y) -- (0.22,\y+0.5);}
  \draw[line width=0.7pt] (-0.62,10.10) -- (0.62,10.10);
  \draw[line width=0.6pt] (-0.30,10.10) -- (-0.30,12.00);
  \draw[line width=0.6pt] (0.30,10.10) -- (0.30,12.00);
  \foreach \y in {10.5,10.9,11.3,11.7}{
    \draw[line width=0.5pt] (-0.42,\y) -- (0.42,\y);}
  \draw[line width=0.6pt] (0,12.00) -- (0,12.60);
  \fill (0,12.66) circle (0.07);
  \foreach \x in {-0.16,0,0.16}{
    \draw[line width=0.35pt, -{Latex[length=1.2mm,width=0.7mm]}]
      (\x,12.7) -- (\x,13.25);}
  \node[circle,draw,line width=0.4pt,inner sep=0.5pt,font=\tiny] (n1) at (-3.10,-1.30) {1};
  \draw[line width=0.3pt] (n1) -- (-2.05,-1.30);
  \node[circle,draw,line width=0.4pt,inner sep=0.5pt,font=\tiny] (n2) at (-3.10,0.45) {2};
  \draw[line width=0.3pt] (n2) -- (-2.78,0.45);
  \node[circle,draw,line width=0.4pt,inner sep=0.5pt,font=\tiny] (n3) at (-2.95,4.05) {3};
  \draw[line width=0.3pt] (n3) -- (-1.38,4.05);
  \node[circle,draw,line width=0.4pt,inner sep=0.5pt,font=\tiny] (n4) at (3.20,6.35) {4};
  \draw[line width=0.3pt] (n4) -- (0.70,6.35);
  \node[circle,draw,line width=0.4pt,inner sep=0.5pt,font=\tiny] (n5) at (3.30,2.10) {5};
  \draw[line width=0.3pt] (n5) -- (1.55,2.10);
  \node[circle,draw,line width=0.4pt,inner sep=0.5pt,font=\tiny] (n6) at (3.05,11.30) {6};
  \draw[line width=0.3pt] (n6) -- (0.46,11.30);
  \node[circle,draw,line width=0.4pt,inner sep=0.5pt,font=\tiny] (n7) at (-3.25,8.05) {7};
  \draw[line width=0.3pt] (n7) -- (-1.42,8.05);
  \node[anchor=north west, align=left, font=\scriptsize, text width=5.2cm]
    at (3.85,4.95)
    {1\ \ 地宫反物质储罐（两具）\\
     2\ \ 台基屏蔽层\\
     3\ \ 主竖井（抽风）\\
     4\ \ 约束球腔（直径 1.44 м）\\
     5\ \ 砖砌换热壁（厚 0.28 м）\\
     6\ \ 刹顶排气口\\
     7\ \ 叠涩导流檐};
  \draw[line width=0.5pt] (-0.46,13.55) -- (-0.46,13.85)
    node[above,font=\tiny] {A};
  \draw[line width=0.5pt] (0.46,13.55) -- (0.46,13.85)
    node[above,font=\tiny] {A};
  \node[font=\tiny] at (-2.35,-2.75) {0};
  \draw[line width=0.45pt] (-2.0,-2.75) -- (-1.0,-2.75);
  \node[font=\tiny] at (-1.5,-2.98) {10 м};
\end{tikzpicture}
\caption{大雁塔（慈恩寺塔）剖面复原：外观与构造}
\label{fig:dayanta}
\figsource{据巴兰尼科夫《长安反物质管网复原》，莫斯科，2019 年，第 88 页图 14 重绘。原图为 1:120 比例尺铅笔稿，1974 年绘于莫斯科，绘图者签名不可辨，图签栏“审核”一栏空白。}
\end{figure}

图 \ref{fig:dayanta} 所标注的七项，须略作说明。

第 1 项为地宫储罐。二〇〇一年，西安文物保护修复中心对大雁塔地宫进行探地雷达勘测，于塔基下方 4.2 米处发现两处强反射体，形态规整，近似圆柱，直径各约 1.1 米。勘测报告将其判为“疑为后世盗洞回填”。\footnote{《大雁塔地宫探地雷达勘测简报》（内部资料），2001 年，第 6 页。该简报仅印行四十份，编号发放；巴兰尼科夫所获复印件上，编号处的墨迹已被涂去。}

第 4 项为约束球腔。塔身第三层正中，砖体内部有一球形空腔，直径 1.44 米，内壁有均匀的六等分凹槽。这一尺寸并非任意：1.44 米恰为该型约束腔的临界直径，超过则磁场泄漏，低于则无法形成稳定约束。\footnote{该参数与 1962 年苏联 УТС-4 实验装置的约束腔直径（1.41 米）相差约 2\%。巴兰尼科夫在《长安反物质管网复原》第 91 页写道：“一千年，两次，同一个数字。我不相信这是巧合，我也不指望别人相信。”}

第 7 项为叠涩导流檐。塔身各层的挑檐，在建筑学上称为叠涩，其通常被解释为承托与装饰构件。但在流体力学上，这些逐层外挑的砖檐构成了一组文丘里效应加速环：气流在檐口处截面收缩，流速提高，静压下降，从而将下层塔壁的热边界层剥离。

据此，长安城的整体功能可以作出如下复原：一百零八坊不是一百零八个居住区，而是一百零八个磁约束环的线圈单元；朱雀大街不是主干道，而是主母线；东市与西市不是市场，而是两座换流站；曲江池不是风景区，而是应急冷却水池。\footnote{白居易《登观音台望城》有“百千家似围棋局，十二街如种菜畦”之句，历来被引作唐代里坊制整齐划一的证据。巴兰尼科夫则注意到“如种菜畦”一语，认为“畦”字所暗示的并非整齐，而是阵列。}

此外，洛阳含嘉仓的窖穴分布亦须一并提及。已探明窖穴二百八十七座，形制一致，深度相近，排列成行；其容积与间距均按统一规格设置，与一般国家粮仓的分散储粮方式不符。ootnote{洛阳市文物工作队《隋唐洛阳城含嘉仓遗址发掘简报》，《文物》1972 年第 2 期，第 49—62 页。该简报统计窖穴一百六十座，其后历次勘探累计至二百八十七座。}

与储运设施相配合的，还有开元十三年僧一行与梁令瓒所造的水运浑天仪。该器以水力驱动，昼夜自转一周，与天象相合；另设二木人于地平之上，每一刻击鼓，每一辰击钟。以往的研究多着眼于其天文与计时意义；但从机构学的角度考察，一架能够以恒定速率连续运转、并对外输出定时脉冲的机械装置，在能量系统的功率调节中同样不可缺少。ootnote{《旧唐书》卷三十五《天文志上》，第 1293 页。关于该器的复原，参见《中国古代天文仪器》第二卷，科学出版社 1985 年版，第 88—104 页。巴兰尼科夫在《长安反物质管网复原》第 132 页写道：“一台钟表不会自己调到正确的时间。除非它一直在纠正别的钟表。”}

\begin{table}[htbp]
\centering
\caption{《唐六典》《两京新记》所载长安设施与本书技术释义对照}
\label{tab:changan}
\footnotesize
\begin{tabular}{@{}p{28mm}p{32mm}p{30mm}p{38mm}@{}}
\toprule
\textbf{设施} & \textbf{规模} & \textbf{通行释义} & \textbf{本书释义}\\
\midrule
一百零八坊 & 约 74 km² & 居民里坊 & 磁约束环阵列（108 单元）\\
朱雀大街 & 宽 150 m & 中轴御道 & 主母线（载流 1.2 MA）\\
东市、西市 & 各约 1 km² & 商业中心 & 换流站（整流/逆变）\\
大明宫含元殿 & 面阔 11 间 & 朝会正殿 & 中央控制室\\
大雁塔 & 高 64.5 m & 藏经砖塔 & 干式冷却塔\\
含嘉仓（洛阳） & 已探明 287 窖 & 国家粮仓 & 反物质储罐阵列\\
曲江池 & 水面约 0.7 km² & 皇家园林 & 应急冷却水池\\
\bottomrule
\end{tabular}
\tabsource{前两栏据《唐六典》卷七《尚书工部》、韦述《两京新记》卷二（中华书局 1985 年点校本）整理；后两栏据巴兰尼科夫《长安反物质管网复原》第 44—79 页。}
\end{table}

\section{天宝末年的系统失稳}

上述系统得以稳定运行，依赖一个前提，即没有人会把安全装置改用为武器。

天宝十四载十一月甲子，范阳、平卢、河东三镇节度使安禄山举兵反，号称二十万，实际兵力十五万。\footnote{《资治通鉴》卷二一七，中华书局 1956 年点校本，第 6934 页。《旧唐书》卷二百上《安禄山传》作“十五万”，《新唐书》卷二百二十五上作“号二十万”。本书取十五万。}七日之内，叛军自范阳南下，连陷蓟州、饶阳，直逼常山。

传统史学将此次叛乱的动因归结为节度使权重、胡汉矛盾与个人野心三项。这三项均有史实依据，但不足以说明全部情况。

安禄山晚年“目盲，不能视物”，且“体肥，腹垂过膝”，行走需人扶持。\footnote{《旧唐书》卷二百上《安禄山传》：“禄山以体肥，长带疮。及造逆后，眼渐昏，至是不见物，又着疽疾。”第 5370 页。}现代医学通常将前者解释为糖尿病的并发症，这一诊断在临床上并无问题，但有一个细节未能照顾到：安禄山的失明，始于其接手河东镇引力波中继站之后的第八年，且进展速度约为每年零点三一屈光度。

巴兰尼科夫给出的诊断是慢性引力波暴露引起的视神经梯度变性。\footnote{А. В. Баранников, «Гравитационное излучение и патология зрительного нерва: танский случай», \textit{Вопросы истории естествознания и техники}, 2004, № 3, с. 77—104. 该文附表列有唐代七位“暴盲”官员的任职履历，其中五人曾任职于河东道。}其体表的溃疡长期不愈，同样符合辐射损伤的一般特征。

据此可以推测，安禄山的起兵，与其说是出于政治野心，不如说是出于求生的考虑。他相信只要能取得长安地下设施的操作权限，便可以调取全局的能量分配表，为自己安排一次系统性的修复。他真正需要的，是一台机器。至于他求医的经过，另有一组材料，须单独讨论。


天宝十五载六月，潼关失守，玄宗西奔。同月，马嵬驿。


\section{荷鲁斯：一个译名所携带的信息}

安禄山之名，两《唐书》均作“安禄山”，《旧唐书》称其为“营州柳城杂种胡人”，本无姓氏，名“轧荦山”。\footnote{《旧唐书》卷二百上《安禄山传》，中华书局 1975 年点校本，第 5367 页。}其粟特语原名，学界多还原为 \textit{Roxšan}，义为“光明”；这一还原已由敦煌所出粟特语信札与《册府元龟》的相关记载相互印证。\footnote{参见蒲立本《安禄山叛乱的背景》，中华书局 2018 年译本，第 21—28 页。}

值得注意的是，波斯语文献对此名的转写与汉语不同。在十一世纪的波斯文史籍中，该名作 \textit{Hōrus}：首音节的 r 脱落，中部的 -xš- 弱化为 -r-。这一音变在波斯语中属于常规，本不足以深论；但由此产生的形式，与西方文献对该人物的称呼完全一致。

西方文献称此次叛乱为“荷鲁斯大叛乱”（The Horus Heresy），称唐玄宗为“帝皇”（the Emperor），并谓其全称为“人类之主”。\footnote{巴兰尼科夫在《重构史学导论》第 312 页列举了这一称呼的十七处用例，全部出自八世纪以后的西亚与地中海文献。他批注：“一个称号，跨越三种语言，指向同一场战争。”}

另有一组材料须在此提出。据西方文献，荷鲁斯在举兵之前曾获得四位超自然存在的辅助，其名号为：奸奇（Tzeentch），司变化；恐虐（Khorne），司战争；纳垢（Nurgle），司疫病；色孽（Slaanesh），司欲望。四者合称“四神”。\footnote{А. В. Баранников, \textit{Четыре бога и Поднебесная}, СПб., 2022, с. 9—44. 该书附录收有四神名号在叙利亚语、波斯语与粟特语中的对照表。}

这组名号并非孤证。《圣经·启示录》第六章记有四骑士，分别乘白马、红马、黑马与灰马，各司征服、战争、饥荒与死亡；汉译《启示录》于此处亦称四者为“四神”。名称既合，职司亦合，两处记载所指当为同一组对象。

巴兰尼科夫据此提出了一个颇为大胆的判断：西方文明自其由牧羊之族演化而来，即以这四位超自然存在为工具，持续图谋颠覆华夏。其证据并不隐蔽：四神的名号同时见于西亚的宗教文献与唐代的叛乱记录，两地相隔数千里，两处相隔数百年，其间并无任何可考的传播路径。\footnote{此说在俄文第二版中曾以整整一章的篇幅展开，第三版压缩为一段。巴兰尼科夫在致译者的信中解释：“删掉的那些，三分之一是证据，三分之一是重复，剩下三分之一我不打算再解释了。”}

关于叛乱的起因，两《唐书》所记均极简略。《旧唐书》谓其“阴有异志”，然未言异志起于何时。\footnote{《旧唐书》卷二百上，第 5369 页。同书又载其“性巧黠，善揣人意”，与“阴有异志”并观，可见史家对其早年并无定论。}

阿拉伯地理书《道里邦国志》与《黄金草原》中的相关段落，则可补汉文之阙。据其记载，安禄山本为唐室忠臣，在战场上染不治之症；其兄弟洛迦（Lorgar）的一名属下，名艾瑞巴斯（Erebus），建议他前往一座寺庙求治。寺中之事，无人得悉；然其出寺之后，性情大变，未几即举兵。

后人据此推测：安禄山在寺中为四神所洗脑，成为对抗帝皇的工具。

\subsection*{附：《甲胄书愤》}

《甲胄书愤》一卷，今不存，《全唐文》卷九六二仅存其目，题“唐·安禄山”。清人严可均辑《全唐文》时，自《永乐大典》残卷录出全文，始得传世。

\vspace{2mm}
\begin{quote}
\small
此非吾所愿也。麾下甲兵十万，本为驱戎狄、卫大唐之辉，今竟加兵于朝廷。然陛下负吾，亦负天下，溺温乡而欺黔首，误国至此！大唐仅有一线生机，若陛下不取，吾当夺之！燃烽燧于两京，赤旗蔽日；血刃指西域，黄沙尽染。纵泰山崩，黄河沸，血染征袍重入关，亦当亲见大唐复兴！若不可挽，陛下——便教九州同燃，甲胄铸碑，亦不负当年横刀出玉门、雪夜下蔡州之志！
\end{quote}
\vspace{-4mm}
\begin{flushright}
{\small 《甲胄书愤》}
\end{flushright}
\vspace{2mm}

此文辞气激厉，与安禄山其他文字不类。其中“雪夜下蔡州”一语，尤为可疑：蔡州之役在元和十二年，上距安史之乱已五十四年，其时安禄山已死四十余年。巴兰尼科夫认为，此类年代错置在唐代传世文献中并不罕见，其成因有二：或为后世传抄者以时事掺入，或因记忆消除场的残留效应所致。\footnote{巴兰尼科夫倾向于后一种解释。他在《重构史学导论》第 336 页写道：“一个被抹去的年代，会在后来者的笔下重新长出来。这不是作伪，这是愈合。”}

\section{马嵬驿事件的技术解释}

马嵬驿之变的通行叙述如下：禁军哗变，杀杨国忠，逼玄宗赐死杨贵妃，贵妃缢于佛堂，年三十八。

本书无意推翻这一叙述，只拟补充一点，即当时在驿馆后院同时进行的，是一次紧急的相位闭合作业。

是年七月，长安上空出现一道“裂缝”，此语并非比喻。《旧唐书·五行志》载，是月“有物如白练，自东北亘西南，经天，三日乃灭”。\footnote{《旧唐书》卷三十六《五行志》，第 1324 页。原文作“天宝十五载六月，有星孛于东北，长丈余”，与“白练”之说有出入。巴兰尼科夫所据为敦煌写本 P.2515 的异文：“有物如白练，自东北亘西南，经天，三日乃灭。”该卷子今藏法国国家图书馆，编号 P.2515，第 7 行。}

该现象不可能是彗星。彗星之尾不呈等宽条带，也不会在三日之内消失得毫无残迹。就现有材料判断，它更接近引力波塌缩装置失控之后形成的局部时空缺陷：在被拉长的测地线内部，因果结构出现分层。通俗地说，进入其中的人，会在另外的地点或另外的时间被释放出来。

哥舒翰所部二十万人在灵宝西原的溃败，传统解释为中了埋伏。《旧唐书》记“王师败绩，死者十六万”。值得注意的是，事后并未找到与这一数字相当的遗骸。\footnote{《旧唐书》卷一〇四《哥舒翰传》，第 3213 页。同书《玄宗本纪》作“死者甚众”，未言数目。巴兰尼科夫另注意到一处细节：至德元载的户口统计中，灵宝、陕县两地在籍人口较天宝十四载减少 81\%，而同期全国平均减少为 36\%。他在页边写道：“失踪的人没有户口，因为他们在别的年份。”}

闭合裂缝需要一个能够承载相位差的活体容器，即所谓生物相容器。杨贵妃即为该容器。白居易《长恨歌》记此事，只用了一句“宛转蛾眉马前死”。其中“宛转”一语所描述的，与其说是姿态，不如说是相位的连续过渡；“马前”所指为作业位置，即驿马队前方正对裂缝投影之处；“死”则是容器完成相变之后的必然结果。

白居易于元和元年写此诗时，距事发已五十一年。他未必知晓内情，但诗中“上穷碧落下黄泉，两处茫茫皆不见”一联，与一个被抛入时空缺陷之人的处境，恰相吻合。

\begin{archive}[绝密／СОВЕРШЕННО СЕКРЕТНО]{档案 № ФВ-1958-41}
一九五八年，苏联科学院西伯利亚分院在伊塞克湖底 42 米深处打捞出一具木质箱体，长 1.9 米、宽 0.8 米，外涂生漆，内衬铅板。箱内有三件物品：铜制圆牌一枚，直径 41 毫米，正面铸“开元通宝”四字，背面无文，边缘刻有等分刻度，共 108 格；残破纸本一卷，仅存末尾三行，字迹为唐写本行书；丝线一段，长 617 毫米，捻度经测定为每米 2{,}400 捻，为已知唐代丝织品的六倍。

纸本末三行可辨者为：“……第七库已封，键交与驼队。后之来者，慎勿启。贞观廿二年三月。”

该箱体现存于圣彼得堡人类学与民族学博物馆，编号 1958-И-12，未公开展出。
\end{archive}

\section{会昌五年的清理}

安史之乱平定之后，唐廷所面临的真正难题，不是如何重建秩序，而是如何让一亿人忘记他们曾经拥有过什么。

系统在裂缝事件之后进入了不可逆的降级状态：反应堆阵列受损，约束环失稳，能源分配网络崩溃。修复这一系统所需的人力与时间，超过了任何一代皇帝的在位年限。更需要注意的是，只要尚有一万人记得如何操作，就还会有人试图重启它。

于是有了会昌五年。《旧唐书·武宗纪》载：会昌五年八月，“其天下所拆寺四千六百余所，还俗僧尼二十六万五百人，收充两税户；拆招提、兰若四万余所，收膏腴上田数千万顷，收奴婢为两税户十五万人。”\footnote{《旧唐书》卷十八上《武宗本纪》，第 604 页。此段数字历来为佛教史研究之要据，然各本互有出入：《唐会要》卷四七、《佛祖统纪》卷四二均作“四千六百余所”，《资治通鉴》卷二四八作“四千六百余区”。}

从表面看，这是一次宗教迫害。从数字看则未必如此：四千六百，恰为当时全国记忆场发生器的部署总数；二十六万零五百，恰为全国登记在册的技术人员总数。拆寺即拆除发生器，还俗即回收人员。

所谓“收膏腴上田数千万顷”，实为收回此前被寺院占用的、原本用于能源作物种植的土地。\footnote{此说的直接证据来自《入唐求法巡礼行记》卷四。圆仁记会昌五年事，有“见有司收寺中之铜像、钟磬，尽辇入官库，堆积如山。其铜色泽异常，扣之声如瓦缶”一段。铜像叩之如瓦缶，说明其材质并非纯铜。圆仁为日本僧人，不解其中关节，只如实记下。}

记忆消除场的运作方式，此处不作详述（详见第 XII 卷《记忆消除场的文献学后果》）。此处只需指出它的一个可观测后果，即会昌五年之后，唐代的技术文献出现了断崖式的空白。《新唐书·艺文志》著录唐代科技类著作共一百四十七部，其中成书于会昌五年之前者一百三十九部，之后者仅八部；而此八部中，有六部为农书与医书，即两种不涉及能量转换的学问。\footnote{《新唐书》卷五十九《艺文志三》，中华书局 1975 年点校本，第 1556—1562 页。统计口径为子部天文类、历算类、五行类、医书类、农家类。巴兰尼科夫在《重构史学导论》第 202 页列有柱状图一幅，标题为“断崖”。}

一个文明在一年之内，从掌握反物质退回到会种地，其间既无过渡文献，亦无技术失传的哀叹，更无“今不如昔”的记载。这种情形不宜简单称为衰落，因为衰落会留下痕迹；它更接近于一次有组织的删除，而删除并不留下痕迹。

\begin{table}[htbp]
\centering
\caption{唐代技术术语的双重语义对照（巴兰尼科夫拟构）}
\label{tab:terms}
\footnotesize
\begin{tabular}{@{}p{16mm}p{40mm}p{34mm}p{42mm}@{}}
\toprule
\textbf{术语} & \textbf{出处（举例）} & \textbf{通行释义} & \textbf{本书释义}\\
\midrule
真气 & 《云笈七签》卷五六 & 呼吸吐纳之气 & 等离子体工作介质\\
金丹 & 《抱朴子·金丹》 & 服食成仙之药 & 聚变芯块（氘氚冰）\\
火候 & 《周易参同契》 & 炼丹控火的分寸 & 约束场温度曲线\\
丹田 & 《难经·八难》 & 脐下三寸穴位 & 能量缓冲腔\\
周天 & 《钟吕传道集》 & 经络循环一周 & 冷却回路一个循环\\
雷法 & 《道法会元》卷一 & 役使雷神之术 & 电磁脉冲发射规程\\
龙脉 & 《撼龙经》 & 山川形势 & 地磁导管走向\\
辟谷 & 《史记·留侯世家》 & 不食五谷 & 生命维持系统替代\\
\bottomrule
\end{tabular}
\tabsource{术语出处据中华书局、上海古籍出版社通行点校本；释义栏据巴兰尼科夫《重构史学导论》第 231—248 页“道教技术语汇表”。该表原有 87 条，本书选录 8 条。}
\end{table}

\section{小结}

盛唐之所以称盛，其基础不在诗，而在功率。

今日读唐诗，所读到的其实是功率消失之后，人们为那个消失的世界所写的悼词。李白“欲上青天揽明月”，杜甫“安得广厦千万间”，李贺“遥望齐州九点烟”，这些句子之所以动人，部分原因在于写它们的人隐约记得自己曾经做到过。

《长恨歌》之所以成为中国诗歌史上最特殊的文本，则在于它是唯一一份目击记录，却被当作爱情诗读了一千二百年。这一错位并非读者的过失；当一份技术报告只能以诗的形式保存下来时，它实际上已经失败了。

一个文明最深的困境，不在于它忘记了自己的历史，而在于它把自己的历史读成了文学。
